%! $n$th Roots \& Rational Exponents — Unit 6, Lesson 6.1 — Warm-Up
\documentclass[10pt]{article}
\usepackage{algebra2-article}
\usepackage{algebra2-boxes}
\newcommand{\TallMath}[1]{$\displaystyle #1\rule[-1.4em]{0pt}{3.2em}$}

\begin{document}

\pageheader{Unit 6, Lesson 6.1}{Warm-Up}
\namedateperiod

\vspace{0.10in}

\begin{notesbox}{Getting Ready: Skills We'll Use Today}
\small These three skills power today's lesson on $n$th roots and rational exponents. Work quickly.

\vspace{4pt}
\textbf{1. Powers you should recognize on sight.} Evaluate each.
\begin{center}
\renewcommand{\arraystretch}{1.5}
\begin{tabular}{@{}l l l l@{}}
(a) \; $2^5=$ \blank{1.6cm} & (b) \; $3^4=$ \blank{1.6cm} &
(c) \; $(-3)^3=$ \blank{1.6cm} & (d) \; $(-2)^4=$ \blank{1.6cm} \\
\end{tabular}
\end{center}
Two of these bases were negative, but only one \emph{answer} came out negative. Which exponents
keep a negative sign, and which destroy it? \writeline

\vspace{4pt}
\textbf{2. The exponent properties, with whole-number exponents.} Simplify. Leave your answer in
exponent form.
\begin{center}
\renewcommand{\arraystretch}{1.5}
\begin{tabular}{@{}l l l l@{}}
(a) \; $x^3\cdot x^5=$ \blank{1.8cm} & (b) \; $\dfrac{x^7}{x^2}=$ \blank{1.8cm} &
(c) \; $\left(x^4\right)^3=$ \blank{1.8cm} & (d) \; $x^{-2}=$ \blank{1.8cm} \\
\end{tabular}
\end{center}
In (a) you \blank{2.2cm} the exponents; in (c) you \blank{2.2cm} them.

\vspace{4pt}
\textbf{3. Fraction arithmetic --- you will need it in the exponents themselves.} Simplify.
\begin{center}
\renewcommand{\arraystretch}{1.6}
\begin{tabular}{@{}l l l@{}}
(a) \; $\dfrac{1}{2}+\dfrac{1}{3}=$ \blank{1.8cm} &
(b) \; $\dfrac{3}{4}-\dfrac{1}{2}=$ \blank{1.8cm} &
(c) \; $\dfrac{2}{3}\cdot 3=$ \blank{1.8cm} \\
\end{tabular}
\end{center}

\end{notesbox}

\end{document}
